\section{Overview}   


With v3, the personal rubric becomes more focused on:  
\begin{itemize}  
	\item A new, more conscious state: \textbf{Being bored} – as a way to reset myself.  
	\item Identifying and initiating \textbf{catalytic events} that \textbf{perturb} my current environment.  
\end{itemize}  

Previously, the state \textit{Challenged} has been transformed into \textbf{Flow State}. This shift occurred to emphasize the feeling state rather than setting specific challenges—at least, not at this stage.\footnote{  
	Certain aspects of the \textit{Serenity} state fulfill this role as well, particularly with \gls{p_NFM} and Franziska.  
}  

The \textit{flow state} is placed at the bottom because it requires less focus, as there is a natural gravitation toward it.  

On the other hand, \textit{Serenity} remains the key state to prioritize. While it is still central, it does not require as much focus at the moment. With the added responsibility of \gls{p_NFM}, the \textit{reset button} has become more important—where previously, it was not a priority.  

The aspect of \textit{being challenged} is still relevant but is now integrated differently within this framework.  
---
\begin{comment}
	Flow State
	
	There is a link to pertubation.
	- Entropy Task are creeping up all the time. There are mimicaing sometimes the emotional state, in which the challenges in the flow state are achived.
	
	- Restrict Entropy Task -> Small tousend cuts -> Focus on something bigger, requires 
	-- to restrict the small entropy taks - to keep up a organisation | That is something I emotionally are able to resit, but still fall short.
	-- the harter part is, to ignore the small challenges, which taks time away from the bigger one, or what is harter, to ignore them inthe first place, becuase if they are satifieing the need
	--- to fully my diserer: To help close college, so that on the floor, I get greated and respected. 
	
	---That is my main objective: To be have colleges on the floor to great me, need me and respect my work.
\end{comment}

 
 \section{Flow State}
 
 \subsection{History - Shift from Challenge to Flow}
 The Flow State represents the domain of \textbf{deep immersion, creative intensity, and self-forgetting engagement}. It replaces the former category of \textit{Challenged}, not because challenge has disappeared, but because the emphasis has shifted: the goal is no longer to manufacture difficulty — it is to enter a qualitative psychological state.\\
 
 The distinction is crucial.  
 \textit{Challenge} is external and measurable.  
 \textbf{Flow} is internal and experiential.
 
 The structural insight of v3 is that I do not need more challenges — I need the right conditions under which challenge naturally transforms into absorption.
 
 The Flow State is therefore positioned at the bottom of the Decision and Focus Space: it requires the least conscious forcing. When the right ingredients are present, I gravitate toward it automatically. It is not scarce. It is conditional.
 
 The task is not to chase Flow.  
 The task is to design environments that allow it to emerge.
 
 \subsection{General Conditions — Play, Experimentation, Construction}
 
 Flow emerges when three elements converge:
 
 \begin{enumerate}
 	\item A clearly defined but open-ended goal.
 	\item Time and cognitive space to think deeply.
 	\item The ability to test, iterate, and observe consequences.
 \end{enumerate}
 
 Historically, I enter Flow when I:
 \begin{itemize}
 	\item immerse myself in a new environment,
 	\item identify a meaningful problem or direction,
 	\item design a solution (conceptual or technical),
 	\item and test whether reality behaves as I predicted.
 \end{itemize}
 
 What matters is the loop:
 
 \begin{center}
 	Think $\rightarrow$ Construct $\rightarrow$ Test $\rightarrow$ Observe $\rightarrow$ Refine
 \end{center}
 
 This loop transforms abstract thinking into embodied understanding.  
 Flow appears when I am simultaneously playing and building.
 
 \subsection{Integration with Serenity}
 
 Maximizing Flow alone — especially when combined with a strong future orientation — risks over-indexing on stimulation and achievement. However, my long-term objective is not stimulation. It is:
 \begin{itemize}
 	\item A wonderful family,
 	\item and great friendships.
 \end{itemize}
 
 Flow must therefore be integrated, not maximized.
 
 The deeper challenge is not to abandon ambition in favor of Serenity — nor to sacrifice Serenity for intensity.
 
 The challenge is to design a life in which:
 
 \begin{center}
 	Flow fuels growth, \\
 	and Serenity protects meaning.
 \end{center}
 
 If this integration succeeds, challenge no longer competes with well-being — it reinforces it.
 
 
 \subsection{Business Environment}
 
 \subsubsection{Beyond Job Descriptions - Anchor Person}
 \begin{center}
 	This [paragraph] suggests that my ideal professional environment is not defined by title, hierarchy, or function — but by the presence of a \textbf{respected anchor person} and \textbf{a mission} that activates both vectors simultaneously.
 \end{center}
 
 The original question underlying this section was: \textit{What is my job description?}\\
 
 The emerging conclusion is that the concept of a fixed job description is not particularly useful for me. My role adapts depending on the person I work with $\bullet$ the mission being pursued $\bullet$ and the learning opportunity available.
 
 My day-to-day drive is not primarily task-based — it is relational and aspirational, eventhough it looks like solving problems and finishing tasks is what I do. However the \textbf{right environment} is set for this by a relational and apsirational binding to a person, that  I:
 \begin{enumerate}
 	\item believe - and for a while choose to belive -  in a person — their abilities, character, or mission,
 	\item feel that I can genuinely learn from them,
 	\item and sense that my accumulated skills and experience meaningfully contribute to their vision diretly or indirectly.
 \end{enumerate}
 
 Given the two enviroment boundries: 
 \begin{enumerate}
 	\item Upward vector: I grow through exposure to someone I respect.
 	\item Outward vector: I contribute something valuable to their mission directly or people which are of interesst to the anchore person.
 \end{enumerate}
 
 Flow arises in the tension between learning and contributing. If either component is missing, engagement decreases: \textit{If I cannot learn, stagnation appears.} $\bullet$ \textit{If I cannot contribute meaningfully, motivation fades.}\\
 
 \subsubsection{Open Hypothesis — Individual vs. Collective Activation}
 
 An open question remains:
 
 Do I respond emotionally in the same way if a group elects me to achieve something for them, rather than aligning under a person I learn from?
 
 Preliminary intuition suggests a difference.
 
 When reporting to or working closely with a person I respect:
 \begin{itemize}
 	\item the learning component is strong,
 	\item the accountability is personal,
 	\item the emotional activation is immediate.
 \end{itemize}
 
 In a purely group-based mandate:
 \begin{itemize}
 	\item contribution may remain strong,
 	\item but the upward learning vector may weaken,
 	\item and with it, the intensity of Flow.
 \end{itemize}
 
 This distinction is not yet fully validated, but it indicates that:
 
 \begin{quote}
 	Flow, for me, is relationally amplified.
 \end{quote}
 
 
 \section{Pertubation}
 %TODO: From Remindes Update or integrate
 % - Segementie Legung (Jobs)
 % - Solve a problem, move to the next
 % - Long Termin view?!
 
 \subsection{Next Jump: Busniess Enviroment}
 
 \begin{itemize}
 	\item A culture differnet setup. Where other setup of values and qualities exist.
 	\item Based on my flow state requirmentes. I still need to find a person that I'm can learn form, and that gives me a mission to follow through.
 	\item I'm not sure: The team persective might be a differnet focus point onto it self.
 	\item A team can be a team, that designs someting new and has to build a \textit{new world}.  Or a team has the function of providing a continuous service.
 \end{itemize}
 
 
 
 \subsection{Study and Research (not yet Science)}
 
 %TODO Stop here. Above is fine.
 As describt above, the research part is more then learning (study) previous knowledge. It is  a systematics way to answering questions or solving problems.\\ 
 
 From the pertubation point of view: It is important to continuously update the rubric part, in order to align and figure out what is most valuable for one self.\\
 
 The other thing would be, that is is important to \textit{spend time} - as an investment==
 
 \subsubsection{not yet Science: Publish}
 As decribt before.
 
 \begin{comment}
 	
 	-Flow State
 	- I don't have a job deciption. For different people I have different jobs.
 	---: 
 	Normally, I aplease myself to somebody or some group - because
 	--- I believe in them or ther mission 
 	--- and I have the feeling, I can learn something from them.
 	--- and I have the sence, that my inate skills, expericse an ambition for ther vision is of use to them.
 	---
 	
 	Therefore I can tell you who I'm serviing, what there goals are and how I think, I help thenm.
 	--= 
 	
 	-> I can therefore tell you: Maybe where I 
 	
 	
 	
 	This section is simple now.
 	
 	What, does it mean?
 	---
 	It's a partially aims to be in the flow state, and partially to pertibuate myself and invest time in new knowleadge frontier.
 	Currently
 	- IT-Data Science: Where two areas are prevelent: Mathmatics and Computer Science.
 	- Concept of Rubrics
 	
 	Field: Occupaction.
 	
 	
 	
 	Mathematics
 	│
 	├── Foundations
 	│   └── Foundation Understanding
 	│
 	├── Methods
 	│   ├── Statistics
 	│   └── Optimization Algorithms
 	│
 	└── Applications
 	└── Machine Learning
 	
 	Computer Science
 	│
 	├── Data
 	│   └── Data Analysis
 	│
 	├── Systems
 	│   └── Engineering
 	│
 	└── Interface
 	└── Visual Tools
 	
 	
 	Math → gives you the theory
 	CS   → gives you the implementation
 	
 	That is a great cycle: 
 	Mathematics
 	↓
 	Computer Science
 	↓
 	Machine Learning
 	↓
 	Data Systems & Tools
 	↓
 	Product Systems
 	↓
 	Business Outcomes
 	
 	Business Context → Economics → Strategy → Product → Revenue
 	
 	
 	
 	Problem → Economics → Decision → System → ROI
 	
 	Job-Descibtion:
 	- Mathematically-powered software systems that automate economically valuable decisions.
 	- Direct Assitent GL SBM-B
 	
 	Title: 
 	- Decision Systems Engineer
 	- Algorithmic Business Decision Engineer (Assitent SBM-B)
 	---
 	
 	Job Description
 	•	Direct Assistant to the General Leadership (GL) of SBM-B, supporting strategic planning and performance management
 	•	Desiging- mathematically-powered software systems that automate economically valuable decisions.
 	
 	
 \end{comment}
 \section{Perturbation}
 \subsection{General Understanding}
 % Simple -> Brilliant
 % Color
 %\definecolor{Rubric_gray}{RGB}{208,208,208}
 %\definecolor{Rubric_FlowState_Lapis_Lazuli}{RGB}{33,95,154}
 %\definecolor{Rubric_Pertubation_Rust}{RGB}{192,79,21}
 %\definecolor{Rubric_Board_Violet}{RGB}{160,43,147}
 %\definecolor{Rubric_Serentiy_BabyBlue}{RGB}{166,202,236}
 
 Perturbing my current environment is essential for counteracting the decay of value, especially in the flow-state environments I inhabit.
 
 For me, transitioning to new environments often requires a catalytic event rather than a slow, gradual shift. There is usually a barrier to overcome, and the natural gravitas of the flow state makes change difficult. Without a decisive disruption, adaptation tends to stagnate.
 
 \subsection{Horizon: Catalytic Events}
 
 In general, there should be only a few key topics within a given period that drive a push toward a new environment. The five-year horizon provides a better guide to identifying the desired perturbations. This yearly horizon focuses on what you actively commit your mind to.
 
 
 \paragraph{This year: End 2026}
 
 \subsection{Optimization Target: Environment over Income}
 
 To strive for a specific monetary amount is misleading, because income itself is not the primary category to optimize.
 
 The central variable is the \textbf{environment}.
 
 The hierarchy of environmental quality is:
 
 \begin{enumerate}
 	\item An environment with an \textbf{anchor person}, a shared \textbf{mission}, and a the anchor person to deliberately learn from.
 	\item An environment with an \textbf{anchor person}, a shared \textbf{mission}, and a person to deliberately learn from.
 	\item An \textbf{isolated but highly challenging situation}, independent of an anchor person.
 \end{enumerate}
 
 The assumption is that consciously selected and intensified challenges increase the income.\textbf{Competition is a key for negoiation.} 
 
 \paragraph{Constraint 1: Income as Stabilizer}
 
 Income is not the optimization goal, but a constraint variable. It must enable:
 \begin{itemize}
 	\item Stability within the social cycle.
 	\item Fun challenge/ Truth motivation:  There is economic headroom in the range of \textbf{160,000–200,000 €} to better match shared life quality expectations.  _
 \end{itemize}
 
 
 Thus, income serves relational stability — not ego validation.
 
 \paragraph{Constraint 2: Flexibility}
 
 The second constraint is maintaining or regaining sufficient flexibility to continuously pursue the long-term objective:
 
 \begin{center}
 	\textbf{Tolle Familie und super Freunde.}
 \end{center}
 
 Any environment that maximizes growth but destroys flexibility is suboptimal.
 The optimal configuration is therefore:
 
 \begin{center}
 	\textit{High-intensity environment} $\cap$ \textit{Relational stability} $\cap$ \textit{Long-term flexibility}
 \end{center}
 
 % In Reminder: Note there different tactis, how to arhieve it or aproach it.
 
 

 
 \paragraph{Longtermin: 20 Year Horizon Catalytic Events: 20 Years}
 
 The push is that, during these periods, a jump into a new environment is required to sustain long-term growth and prevent decay. These five-year blocks are arbitrary intervals but serve to provide a long-term perspective, helping to clarify where you want to work towards.
 
 
 

 
 %%%%%%
 \textit{I'm came with abilities and an emotional imprint.}
 
 I can build up domain knowdledge an update this, to combine with my concitive abilities. How ever, think about how I feel, how I want to feel and how I may need to change how I feel, in order to satifice my core imprint, this the goal for this - Homostatice
 
 I can't do much about the former, except build better prediction models for my self and a cummilate knowledge - but even this is limited. The later however, is more malible, or at least essesing it.

